\documentclass[10pt]{article}
\usepackage[margin=1in]{geometry}
\usepackage{amsmath,amssymb,amsthm,booktabs,array}
\usepackage[T1]{fontenc}
\usepackage{lmodern}
\usepackage[hidelinks,breaklinks]{hyperref}
\allowdisplaybreaks[1]
\newtheorem{theorem}{Theorem}
\newtheorem{corollary}[theorem]{Corollary}
\newcommand{\Q}{\mathbb{Q}}
\newcommand{\Z}{\mathbb{Z}}
\newcommand{\coeff}[2]{[#1^{#2}]}
\title{Mathar's inverse-quintic recurrence conjecture\\
and extensions of Barry's cubic formulas}
\author{Haobo Ma\\\small ChronoAI Pte Ltd, Singapore
\and Wenlin Zhang\\\small National University of Singapore\\\small ChronoAI Pte Ltd, Singapore}
\date{November 14, 2026}
\begin{document}
\maketitle
\begin{abstract}
We prove R.\ J.\ Mathar's ``Conjecture D-finite with recurrence'' in the third
FORMULA line of OEIS A392976, revision 19 (January 31, 2026; conjecture dated
January 29, 2026), and extend the cubic coefficient identities proved by
Barry (2011, Theorem 3 and Example 7). The coefficients of the formal inverse of
$t((1-t^2)^2-t)$, divided by $x$, satisfy the stated six-term recurrence at
every integer $n\geq5$. A change of variable reduces the assertion to an
explicit polynomial identity. The same argument gives a nonzero homogeneous
differential equation of order five; a trace argument lowers the existence
bound to order four. Lagrange inversion gives an integer-exponent
family containing the entry's binomial-sum formula and its three cross-referenced
sequences, as well as all positive integral powers of their generating functions.
Further corollaries cover the cubic entry's alternate sum, differential recurrence,
and derivative formulas.
\end{abstract}

\section{The conjecture and its exact recurrence}

Let
\[
 f(t)=t((1-t^2)^2-t)=t-t^2-2t^3+t^5.
\]
Since $f(0)=0$ and $f'(0)=1$, there is a unique $U\in x\Q[[x]]$
with $f(U(x))=x$; it also satisfies $U(f(t))=t$. Put
\[
 A(x)=U(x)/x=\sum_{n\geq0}a_nx^n.
\]
This is Seiichi Manyama's sequence A392976. Its third FORMULA line credits
the recurrence conjecture to R.\ J.\ Mathar~\cite{oeis}.
We prove that conjecture below by transporting a differential operator
through the inverse. Barry's earlier coefficient argument for the cubic
inverse~\cite[Theorem 3 and Example 7]{barry} is the $r=1$ case of the
family in Section~3. Schuetz and Whieldon~\cite[Theorem 1.9 and Section 3]{sw}
and Birmajer, Gil and Weiner~\cite[Section 4]{bgw} also prove the
cubic polygon-dissection formulas; none of these results asserts Mathar's
specific quintic recurrence.
Liszewska and M\l{}otkowski~\cite[Proposition 6.1, equation (19)]{lm}
give a further Lagrange-inversion proof of the cubic coefficient formula.
The source's offset is zero. It specifies no recurrence onset: $n\geq5$
is the range in which all six sequence arguments are nonnegative.

Write the six signed coefficient polynomials as follows:
\begin{align*}
c_0(z)={}&283z(z-1)(z-2)(83022305213z-178385058984)(z+1),\\
c_1(z)={}&-6z(z-1)(z-2)\bigl(12196786904129z^2\\*
 &\hspace{36mm}-45110300936020z+52641391340361\bigr),\\
c_2(z)={}&-6(z-1)(z-2)\bigl(67790996421377z^3\\*
 &\hspace{22mm}-248299346995054z^2+202497447227007z\\*
 &\hspace{22mm}+3061418504418\bigr),\\
c_3(z)={}&-2(z-2)\bigl(192167357944360z^4-559010679688491z^3\\*
 &\hspace{22mm}-2120668219560541z^2+9123363136130610z\\*
 &\hspace{22mm}-8342981628089220\bigr),\\
c_4(z)={}&5\bigl(15640122296875z^5-565964885220500z^4\\*
 &\hspace{12mm}+4989857715427325z^3-18667678380008380z^2\\*
 &\hspace{12mm}+32085126156301848z-20925325629575136\bigr),\\
c_5(z)={}&10(5z-17)(5z-18)(5z-21)(5z-19)\\*
 &\hspace{12mm}\cdot(22509687685z-60308240636).
\end{align*}

\begin{theorem}\label{thm:main}
For every integer $n\geq5$,
\begin{equation}\label{eq:rec}
 \sum_{k=0}^{5}c_k(n)a_{n-k}=0.
\end{equation}
The series $A$ is D-finite: it satisfies a homogeneous linear differential
equation with polynomial coefficients, not all zero, of order five.
Together with $a_0,\ldots,a_4=1,1,4,15,67$, equation~\eqref{eq:rec}
uniquely determines the sequence.
\end{theorem}

\section{Proof by a polynomial identity}

Algebraic series are D-finite~\cite[Theorem 2.1]{stanley}, so existence
of some differential equation already follows from $f(U)=x$. The point
here is to verify the six particular polynomials in Mathar's conjecture.
The denominator-clearing method is the inverse-variable Euler form of
the algebraic differentiation procedure described by
Bostan et al.~\cite[Section 2.1]{bostan}.

Set $\theta=x\,d/dx$ and define
\[
 L_A=\sum_{k=0}^{5}x^kc_k(\theta+k),\qquad
 L_U=\sum_{k=0}^{5}x^kc_k(\theta+k-1).
\]
Here $x^k$ multiplies on the left; the polynomial in $\theta$ acts first.
The identity $\theta(xB)=x(\theta+1)B$ implies
\begin{equation}\label{eq:shift}
 L_U(xB)=xL_AB.
\end{equation}
We prove $L_UU=0$ after composing with $f$.

Let $H(t)=f'(t)=1-2t-6t^2+5t^4$, a unit in $\Q[[t]]$. The formal
chain rule gives
\begin{equation}\label{eq:transport}
 (\theta B)(f(t))=\frac{f(t)}{H(t)}\frac{d}{dt}B(f(t)).
\end{equation}
Define polynomials $N_0=t$, $N_1=f$, and, for $1\leq j\leq4$,
\begin{equation}\label{eq:num}
 N_{j+1}=f\bigl(HN_j'-(2j-1)H'N_j\bigr).
\end{equation}
Since $U(f(t))=t$, induction using~\eqref{eq:transport} proves
\[
 (\theta^jU)(f(t))=\frac{N_j(t)}{H(t)^{2j-1}}
 \quad(1\leq j\leq5).
\]
For example $N_1=t-t^2-2t^3+t^5$ and
\[
 N_2=t-3t^2+2t^3+10t^4-9t^5-18t^6-12t^7
       +12t^8+23t^9-3t^{10}-18t^{11}+5t^{13}.
\]
Define integers $w_{kj}$ by
\begin{equation}\label{eq:weights}
 c_k(z+k-1)=\sum_{j=0}^{5}w_{kj}z^j.
\end{equation}
The required finite identity in $\Z[t]$ is
\begin{equation}\label{eq:certificate}
 \sum_{k=0}^{5}f(t)^k
 \left(w_{k0}tH(t)^9+
       \sum_{j=1}^{5}w_{kj}N_j(t)H(t)^{10-2j}\right)=0.
\end{equation}
All its ingredients are specified by the displayed $c_k$, $f$, $H$ and
recursion~\eqref{eq:num}. Expanding them gives the zero polynomial;
\texttt{verify.py} checks every coefficient by integer arithmetic.
Thus~\eqref{eq:certificate} proves
$H^9(L_UU)(f(t))=0$. Since $H$ is a unit, $(L_UU)(f(t))=0$.
Composition with $U$ gives $L_UU=0$. Equation~\eqref{eq:shift} and
$U=xA$ now imply $xL_AA=0$, hence $L_AA=0$.

Since $\theta x^m=mx^m$, the coefficient of $x^n$ in
$x^kc_k(\theta+k)A$ is $c_k(n)a_{n-k}$ whenever $n\geq k$.
Extracting coefficients of $L_AA=0$ proves~\eqref{eq:rec}
for every $n\geq5$, not just for a finite initial segment.

Each $c_k$ has degree five. Expanding powers of $\theta$ in ordinary
derivatives shows that the coefficient of $d^5/dx^5$ in $L_A$ is
\begin{align*}
 x^5\bigl(&23495312375279-73180721424774x
 -406745978528262x^2\\
 &-384334715888720x^3+78200611484375x^4
 +140685548031250x^5\bigr).
\end{align*}
It is nonzero, proving D-finiteness. Finally $c_0(n)>0$ for every $n\geq5$,
so~\eqref{eq:rec} determines each new term from its five predecessors.
The first five coefficients follow by recursive inversion of $f$.
This completes the proof of Theorem~\ref{thm:main}.

\begin{corollary}\label{cor:order}
The series $A$ also satisfies a nonzero homogeneous linear differential
equation over $\Q[x]$ of order at most four.
\end{corollary}
\begin{proof}
The polynomial $f(T)-x$ is irreducible over $\Q(x)$: it is primitive,
and its quotient ring over $\Q[x,T]$ is $\Q[T]$, an integral domain.
Gauss's lemma then applies. Thus $K=\Q(x,U)$ has degree five over
$\Q(x)$. The coefficient of $T^4$ in $f(T)-x$ is zero, so
$\operatorname{Tr}_{K/\Q(x)}U=0$, and the same holds for $A=U/x$.
In characteristic zero, differentiation extends uniquely to every finite
algebraic extension. Consequently it commutes with the embeddings defining
the trace, so every derivative of $A$ has trace zero. Since the trace of
$1$ is $5$, the trace-zero subspace has dimension four. The five elements
$A,A',\ldots,A^{(4)}$ are linearly dependent over $\Q(x)$.
Clearing denominators gives the stated equation. This is an existence
bound, not a claim that the displayed order-five operator has order four.
\end{proof}

\section{Generating functions and neighboring formulas}

The defining equation $f(U)=x$ immediately yields the other generating-function
identity in A392976:
\begin{equation}\label{eq:gf}
 A(x)=\frac{1}{(1-(xA(x))^2)^2-xA(x)}.
\end{equation}
Its binomial formula is a specialization of the following calculation.

\begin{corollary}\label{cor:family}
Let $r$ be any integer, let $U_r$ be the zero-constant inverse of
$f_r(t)=t((1-t^2)^r-t)$, and write $U_r(x)/x=\sum_{n\geq0}a_n^{(r)}x^n$.
For integers $n\geq0$ and $s\geq1$,
\begin{equation}\label{eq:powers}
 [x^n](U_r(x)/x)^s=\frac{s}{n+s}
 \sum_{k=0}^{\lfloor n/2\rfloor}
 \binom{2n+s-2k-1}{n-2k}
 (-1)^k\binom{-r(2n+s-2k)}{k}.
\end{equation}
Here $\binom{b}{k}=b(b-1)\cdots(b-k+1)/k!$ for any integer $b$,
and $\binom{b}{0}=1$. In particular, for $r\geq1$,
\begin{equation}\label{eq:family}
 a_n^{(r)}=\frac{1}{n+1}\sum_{k=0}^{\lfloor n/2\rfloor}
 \binom{2n-2k}{n}\binom{r(2n-2k+1)+k-1}{k}.
\end{equation}
All $a_n^{(r)}$ are integers, and they are strictly positive when $r\geq0$.
Every $U_r/x$ is algebraic and D-finite. The cases $r=1,2,3,-1$ give,
respectively, A001002, A392976, A392977 and A370839.
\end{corollary}

\begin{proof}
Put $\phi_r(t)=((1-t^2)^r-t)^{-1}$. This is a formal series with
constant coefficient one, also when $r<0$. Lagrange inversion
\cite[Theorem 2.1.1, equation (2.1.1)]{gessel} gives
\begin{equation}\label{eq:lagrange}
 a_n^{(r)}=\frac{1}{n+1}[t^n]\phi_r(t)^{n+1}.
\end{equation}
Indeed $U_r=x\phi_r(U_r)$, so the cited formula applies over $\Q[[x]]$
with a unit constant term. One may obtain the same identity by the formal
residue substitution $x=f_r(t)$ and the fact that a Laurent derivative has
zero residue.
The binomial expansion is
\[
 \phi_r(t)^{n+1}=\sum_{j\geq0}\binom{n+j}{j}
 t^j(1-t^2)^{-r(n+1+j)}.
\]
Only $j=n-2k$ contributes to $[t^n]$. Expanding the last factor gives
\eqref{eq:family}. More generally the power form of Lagrange inversion is
$[x^{n+s}]U_r^s=\frac{s}{n+s}[t^n]\phi_r(t)^{n+s}$.
Expanding first in $t(1-t^2)^{-r}$ and then in $t^2$ gives
\eqref{eq:powers}; thus the formula is valid for negative and zero $r$ as
well. Equation~\eqref{eq:family} is exactly the published sum in A392976 and A392977
\cite{oeis,oeis3}. For $r=1$ its summand equals
$\binom{2n-k}{n+k}\binom{n+k}{k}$, the formula in A001002~\cite{oeis1}.
The ring $\Z[[t]]$ contains $(1-t^2)^r$ and its unit reciprocal,
so coefficient recursion in $U_r=x\phi_r(U_r)$ proves integrality for every
integer $r$. For $r\geq1$, $\phi_r$ has nonnegative coefficients, and the
$k=0$ summand of~\eqref{eq:family} is positive. For $r=0$ one obtains
$a_n^{(0)}=\binom{2n}{n}/(n+1)>0$.

To prove the final assertion, clear the power of $1-U_r^2$ from
$f_r(U_r)=x$. This gives a nonzero polynomial relation over $\Q(x)$,
of degree at most $2r+1$ in $U_r$ for $r\geq1$, degree two for $r=0$,
and degree $2|r|+2$ for $r<0$. In characteristic zero the finite extension
$K=\Q(x,U_r)$ is closed under differentiation: differentiating a minimal
polynomial expresses $U_r'$ in $K$, since that polynomial is separable.
All derivatives of $U_r/x$ therefore lie in $K$. More than $[K:\Q(x)]$
of them are linearly dependent over $\Q(x)$; clearing denominators gives
a nonzero homogeneous polynomial-coefficient differential equation.
\end{proof}

\begin{corollary}\label{cor:signed}
For A370839 the coefficient formula is
\[
 a_n^{(-1)}=\frac{1}{n+1}
 \sum_{k=0}^{\lfloor n/2\rfloor}(-1)^k
 \binom{2n-2k}{n-2k}\binom{2n-2k+1}{k}.
\]
The generating functions of A392977 and A370839 satisfy
$A_r=1/((1-(xA_r)^2)^r-xA_r)$ with $r=3,-1$, respectively.
\end{corollary}
\begin{proof}
Set $r=-1$, $s=1$ in~\eqref{eq:powers}. This proves the formula in
A370839~\cite{oeism1}. Substituting $U_r=xA_r$ in $f_r(U_r)=x$
proves both generating-function identities. In particular $a_4^{(-1)}=-5$:
strict positivity for this signed specialization is not asserted in the source.
\end{proof}

\begin{corollary}\label{cor:cubic}
Let $B=U_1/x=\sum_{n\geq0}b_nx^n$, the generating function of A001002.
Then
\begin{align}
 B&=1+xB^2+x^2B^3,\\
 b_n&=\frac1{n+1}\sum_{j=\lceil n/2\rceil}^n
       \binom{n+j}{j}\binom{j}{n-j},\\
 5n(n+1)b_n&=11n(2n-1)b_{n-1}
       +3(3n-2)(3n-4)b_{n-2}\quad(n\geq2).\label{eq:cubicrec}
\end{align}
For $g(x)=x^2+x^3$ and $D=d/dx$ its two derivative representations are
\begin{align}
 B&=1+\frac1x\sum_{m\geq1}\frac{D^{m-1}g(x)^m}{m!},\\
 B&=\exp\left(\sum_{m\geq1}\frac{D^{m-1}(g(x)^m/x)}{m!}\right).
\end{align}
These are identities in $\Q[[x]]$; each coefficient receives finitely many terms.
\end{corollary}
\begin{proof}
The inverse equation is $U_1=x+U_1^2+U_1^3$, giving the first identity.
Expand $(1-t-t^2)^{-n-1}$ as
$\sum_{j\geq0}\binom{n+j}{j}t^j(1+t)^j$ in~\eqref{eq:lagrange}.
This gives Len Smiley's alternate sum in A001002, also proved in
Barry~\cite[Example 7]{barry}. Replacing $j$ by $n-k$ identifies it
term by term with Hanna's sum.

For~\eqref{eq:cubicrec}, use the proof of Theorem~\ref{thm:main} with
$f_1=t-t^2-t^3$, $H_1=1-2t-3t^2$ and the three signed polynomials
$5z(z+1)$, $-11z(2z-1)$, $-3(3z-2)(3z-4)$.
For the shifted operator on $U_1$, the numerators of its first two Euler
iterates are $f_1$ and $f_1(H_1f_1'-H_1'f_1)$.
Multiplication by $H_1^3$ gives $-f_1^2H_1^3$; equivalently the
corresponding operators satisfy $L_{U_1}U_1=-x^2$ and $L_BB=-x$.
The accompanying integer checker verifies the polynomial identity, and
coefficient extraction proves the recurrence for every $n\geq2$.
This reproves the recurrence
attributed to Smiley in A001002~\cite{oeis1}.

Finally introduce a parameter $z$ and solve $V=x+zg(V)$.
The Lagrange--B\"urmann expansion at $x$ is
$F(V)=F(x)+\sum_{m\geq1}z^mD^{m-1}(F'(x)g(x)^m)/m!$
\cite[Section 2.1]{gessel}. Take $F(t)=t$ for the first representation,
and $F(t)=\log t$ for the second after subtracting $\log x$.
All resulting terms in the latter difference have positive $x$-adic
order, so exponentiation and $z=1$ are legitimate. These are Hanna's
two derivative formulas in A001002.
\end{proof}

\begin{corollary}\label{cor:reciprocal}
Put $C=1/B=\sum_{n\geq0}d_nx^n$ and
$\ell_n=[x^n]xB'/B$ for $n\geq1$. Then
\begin{align}
 b_n&=-\sum_{i=1}^n d_i b_{n-i} &&(n\geq1),\\
 b_n&=-\sum_{i=1}^{n+2}d_i d_{n+2-i} &&(n\geq0),\\
 \ell_n&=[t^n](1-t-t^2)^{-n} &&(n\geq1),\\
 b_n&=\frac1{n+1}\binom{2n}{n}
 {}_2F_1\!\left(-\frac n2,\frac{1-n}2;-2n;-4\right) &&(n\geq1).
\end{align}
The hypergeometric expression means the finite sum through
$k=\lfloor n/2\rfloor$ with summand
$(-n/2)_k((1-n)/2)_k(-4)^k/((-2n)_k k!)$;
$(a)_k=a(a+1)\cdots(a+k-1)$ and $(a)_0=1$.
The coefficients $d_n$ form A217596, and the $\ell_n$ form A213684.
\end{corollary}
\begin{proof}
From $BC=1$ one obtains the first convolution. The cubic equation gives
$C=1-xB-x^2B^2$, hence $C^2=C-x-x^2B$. Comparing degree $n+2$
gives the second convolution, since the two endpoint terms of $C^2$
are $2d_{n+2}$. These are the two reciprocal-sequence formulas in
A001002~\cite{oeis1,reciprocal}.

Write $\phi(t)=(1-t-t^2)^{-1}$. The Lagrange--B\"urmann formula gives
$[x^n]\log B=\frac1n[t^{n-1}]\phi'(t)\phi(t)^{n-1}
=\frac1n[t^n]\phi(t)^n$. Multiplication by $n$ proves the third identity,
including its index convention~\cite{logderivative}.
Finally the $r=1$ summand in~\eqref{eq:family} is
$(2n-k)!/(n!k!(n-2k)!)$. Dividing its value by $\binom{2n}{n}$
gives the displayed terminating hypergeometric summand, proving
Luschny's expression in A001002~\cite{oeis1}.
\end{proof}

The antidiagonal assertion in A001002 follows as well: if
$F(x,y)=1+xF(x,y)^2+xyF(x,y)^3$ with constant term one, substitution
$y=x$ gives $F(x,x)=B(x)$ by coefficient recursion. Barry's Bell-matrix
identity $[t^n]\phi(t)^{n+1}=(n+1)[x^{n+1}]U(x)$
\cite[Theorem 3]{barry} is precisely~\eqref{eq:lagrange}, not an
additional unproved transform.

The cubic's real branch near $x=0$ is
\[
 B(x)=\frac{4\sin\!\left(\frac13\arcsin\frac{27x+11}{16}\right)-1}{3x},
\]
with its removable value $B(0)=1$. Indeed, writing
$v=(4\sin\alpha-1)/3$ gives $v-v^2-v^3=(16\sin3\alpha-11)/27$.
At $x=0$ the displayed branch has $v=0$ and derivative one, so formal
inverse uniqueness proves the identity. This agrees with the OEIS formula;
the denominator printed inside the arcsine in Barry's Example 7 is $6$,
which must be $16$ for this inverse.

The exact minimal differential order and the minimal recurrence order
for A392976 are not determined here; Corollary~\ref{cor:order} gives
the stronger differential-order bound four.

\section{Coverage and verification}

\begin{center}
\small
\begin{tabular}{@{}p{0.27\linewidth}p{0.64\linewidth}@{}}
\toprule
Source statement & Mathematical coverage\\
\midrule
A392976, FORMULA 3 & Proved: Theorem~\ref{thm:main}, exact recurrence for $n\geq5$
and D-finiteness.\\
A392976, FORMULAS 1--2 & Proved: equation~\eqref{eq:gf} and
Corollary~\ref{cor:family} with $r=2$: written proofs.\\
A392977, FORMULAS 1--2 & Proved: same generating-function argument and
Corollary~\ref{cor:family} with $r=3$: written proofs.\\
A001002, binomial formula & Reproved; prior proof Barry (2011), Example 7;
Corollary~\ref{cor:family} with $r=1$:
reproof of the published formula.\\
A001002, polynomial identity, recurrence, derivative forms, antidiagonals
& Reproved: Corollary~\ref{cor:cubic} and the following substitution;
attributions to Smiley and Hanna retained.\\
A001002, trigonometric formula & Proved as stated in OEIS; the $6$ in
Barry's Example 7 is corrected to $16$ above.\\
A001002, polygon interpretation & Prior proof cited:
Schuetz--Whieldon (2014), Theorem 1.9; Birmajer--Gil--Weiner (2015), Section 4.\\
Integer $r$, positive powers & Not posed; proved in
Corollary~\ref{cor:family}. Algebraicity, D-finiteness and integrality
hold for every integer $r$; positivity holds for $r\geq0$.\\
Quintic differential order & Not posed; proved: the trace argument in
Corollary~\ref{cor:order} gives order at most four.\\
A001002, reciprocal convolutions, logarithmic derivative, hypergeometric form
& Proved: Corollary~\ref{cor:reciprocal}; published formulas reproved.\\
A001002, remaining identities and interpretations & Not covered:
parking functions, infinite product and asymptotic formula.\\
A370839, FORMULA & Proved: Corollary~\ref{cor:signed},
reproof of the published formula.\\
\bottomrule
\end{tabular}
\end{center}

The \href{https://github.com/the-omega-institute/trureturing/blob/2b5b552e72596710282b713f6587c7020544535d/D5/S1/Recurrence/Algebraic/A392976InverseRecurrence.lean}{pinned Lean file}~\cite{lean} uses the namespace
\begin{quote}\ttfamily
D5.S1.Recurrence.Algebraic.A392976InverseRecurrence
\end{quote}
Its theorem \texttt{result} proves both inverse compositions, uniqueness of the zero-constant right
inverse, $xA=U$, homogeneous D-finiteness, and the exact recurrence for every
natural index $n\geq5$. These indices represent all integers in the theorem's
range. Its definitions of $f$, $U$, $A$ and $a$ agree with those above.
The initial five coefficient values, their determination of later terms,
the trace-zero order bound, all binomial-sum formulas, the integer-parameter family and the cubic
corollaries have written proofs here; they are not additional conjuncts of
\texttt{result}. The declaration \texttt{polynomial\_certificate}
in the same pinned file is the quintic polynomial identity, and
\texttt{normal\_form} identifies an order-five operator with a nonzero
fifth-derivative coefficient. The accompanying integer-polynomial checker verifies
\eqref{eq:certificate}; it also compares coefficient formulas with recursive
reversion on a finite range specified in \texttt{verification.json}.
The universal recurrence follows from the polynomial identity and the
formal-series argument, rather than extrapolation from those comparisons.

\begin{thebibliography}{99}
\bibitem{stanley}
R.\ P.\ Stanley, \emph{Differentiably finite power series},
European J.\ Combin.\ \textbf{1} (1980), 175--188, Theorem 2.1, p.\ 178.
\href{https://math.mit.edu/~rstan/pubs/pubfiles/45.pdf}{Author-hosted full text}.
\bibitem{bostan}
A.\ Bostan, F.\ Chyzak, B.\ Salvy, G.\ Lecerf and \'E.\ Schost,
\emph{Differential equations for algebraic functions}, ISSAC '07 (2007),
25--32, Section 2.1.
\href{https://arxiv.org/abs/cs/0703121v2}{arXiv:cs/0703121v2}.
\bibitem{barry}
P.\ Barry, \emph{On the Central Coefficients of Bell Matrices},
J.\ Integer Seq.\ \textbf{14} (2011), Article 11.4.3,
Theorem 3 and Example 7, pp.\ 3--4 and 6--7.
\href{https://cs.uwaterloo.ca/journals/JIS/VOL14/Barry3/barry132.pdf}{Full text}.
\bibitem{lm}
E.\ Liszewska and W.\ M\l{}otkowski, \emph{Some relatives of the Catalan sequence},
\href{https://arxiv.org/abs/1907.10725v1}{arXiv:1907.10725v1} (2019),
Proposition 6.1, equation (19), and Section 9.3.
\bibitem{reciprocal}
V.\ Kruchinin and contributors, \emph{A217596}, OEIS, revision 63,
July 29, 2025.
\href{https://github.com/oeis/oeisdata/blob/4c6abcdf8b00fd262030066ebfced671457e5b43/seq/A217/A217596.seq}{Pinned entry}.
\bibitem{logderivative}
P.\ D.\ Hanna and contributors, \emph{A213684}, OEIS, revision 46,
November 5, 2025.
\href{https://github.com/oeis/oeisdata/blob/4c6abcdf8b00fd262030066ebfced671457e5b43/seq/A213/A213684.seq}{Pinned entry}.
\bibitem{sw}
A.\ Schuetz and G.\ Whieldon, \emph{Polygonal Dissections and Reversions of Series},
\href{https://arxiv.org/abs/1401.7194v1}{arXiv:1401.7194v1} (2014),
Theorem 1.9 and Section 3, pp.\ 8--10.
\bibitem{bgw}
D.\ Birmajer, J.\ B.\ Gil and M.\ D.\ Weiner,
\emph{Colored partitions of a convex polygon by noncrossing diagonals},
\href{https://arxiv.org/abs/1503.05242v1}{arXiv:1503.05242v1} (2015),
Section 4, p.\ 5.
\bibitem{oeis}
S.\ Manyama, R.\ J.\ Mathar and contributors,
\emph{A392976}, The On-Line Encyclopedia of Integer Sequences,
revision 19, January 31, 2026; Mathar's FORMULA conjecture dated January 29, 2026.
\href{https://github.com/oeis/oeisdata/blob/4c6abcdf8b00fd262030066ebfced671457e5b43/seq/A392/A392976.seq}{Pinned entry, lines 5--9}.
\bibitem{gessel}
I.\ M.\ Gessel, \emph{Lagrange inversion},
\href{https://arxiv.org/abs/1609.05988v1}{arXiv:1609.05988v1} (2016),
Theorem 2.1.1, equation (2.1.1).
\bibitem{oeis3}
S.\ Manyama and contributors, \emph{A392977}, OEIS,
revision 16, January 31, 2026.
\href{https://github.com/oeis/oeisdata/blob/4c6abcdf8b00fd262030066ebfced671457e5b43/seq/A392/A392977.seq}{Pinned entry}.
\bibitem{oeis1}
N.\ J.\ A.\ Sloane and contributors, \emph{A001002}, OEIS,
revision 183, February 11, 2026; binomial formula attributed to P.\ D.\ Hanna,
March 30, 2005.
\href{https://github.com/oeis/oeisdata/blob/4c6abcdf8b00fd262030066ebfced671457e5b43/seq/A001/A001002.seq}{Pinned entry}.
\bibitem{oeism1}
S.\ Manyama, \emph{A370839}, OEIS,
revision 11, March 3, 2024.
\href{https://github.com/oeis/oeisdata/blob/4c6abcdf8b00fd262030066ebfced671457e5b43/seq/A370/A370839.seq}{Pinned entry}.
\bibitem{lean}
\emph{A392976InverseRecurrence.lean}, source revision\\
\texttt{2b5b552e72596710282b713f6587c7020544535d}.
\href{https://github.com/the-omega-institute/trureturing/blob/2b5b552e72596710282b713f6587c7020544535d/D5/S1/Recurrence/Algebraic/A392976InverseRecurrence.lean}{Pinned Lean declarations}.
\end{thebibliography}

\medskip\noindent\footnotesize
\copyright\ 2026 Haobo Ma and Wenlin Zhang.
Licensed under Creative Commons Attribution 4.0 International (CC BY 4.0).
OEIS entry material is attributed to the OEIS Foundation and its named
contributors; the cited entries are available under CC BY-SA 4.0.
\end{document}
